\section*{अध्याय \ref{ch:b2:heart} --- हृदय और हृद्-चक्र}

\begin{solution}{exo:b2:heart:1}
महाशिरा $\to$ दायाँ \omterm{def:b2:heart:anatomy}{अलिंद} $\to$ त्रिवलनी कपाट $\to$ दायाँ
\omterm{def:b2:heart:anatomy}{निलय} $\to$ फुफ्फुसीय कपाट $\to$ फुफ्फुस \omterm{def:b2:blood-circulation:vessels}{धमनी} $\to$ फेफड़े
$\to$ फुफ्फुस \omterm{def:b2:blood-circulation:vessels}{शिराएँ} $\to$ बायाँ \omterm{def:b2:heart:anatomy}{अलिंद} $\to$ द्विवलनी कपाट $\to$ बायाँ
\omterm{def:b2:heart:anatomy}{निलय} $\to$ महाधमनीय कपाट $\to$ महाधमनी।
\end{solution}

\begin{solution}{exo:b2:heart:2}
भरना (अनुशिथिलन): AV कपाट खुले, अर्धचंद्राकार बंद। समआयतनी
संकुचन: चारों बंद। निष्कासन: अर्धचंद्राकार खुले, AV बंद।
समआयतनी शिथिलन: सब बंद। पहली ध्वनि: प्रकुंचन के आरंभ पर बंद
होते AV कपाट; दूसरी: उसके अंत पर बंद होते अर्धचंद्राकार
कपाट।
\end{solution}

\begin{solution}{exo:b2:heart:3}
\omterm{prop:b2:heart:conduction}{साइनु-अलिंद गाँठ} ($t = 0$; \omterm{def:b2:heart:anatomy}{अलिंदों} के विध्रुवित होते ही P तरंग);
\omterm{prop:b2:heart:conduction}{अलिंद-निलय गाँठ} (लगभग \qty{80}{ms} पर पहुँचा, उस आवेग को
\qty{100}{ms} थामे: यानी PR खंड); हिस बंडल और
\omterm{prop:b2:heart:conduction}{पुर्किंजे तंतु} (\qtyrange{180}{220}{ms}; \omterm{def:b2:heart:anatomy}{निलयों} के विध्रुवित
होते ही QRS); और बाद में उनके पुनर्ध्रुवित होते ही T तरंग।
\end{solution}

\begin{solution}{exo:b2:heart:4}
अपने कार्य-परास के भीतर वह \omterm{def:b2:heart:anatomy}{निलय} तब अधिक निकालता है जब वह अधिक
भरा गया हो। यदि दायाँ \omterm{def:b2:heart:anatomy}{निलय} अधिक पंप करे, तो बायाँ कुछ धड़कनों बाद
अधिक पाता है और, उसी नियम से, अधिक पंप करता है: कोई भी
असंतुलन बिना किसी संकेत के स्वयं सुधर जाता है।
\end{solution}

\begin{solution}{exo:b2:heart:5}
\omterm{prop:b2:heart:cycle}{आघात-आयतन} \qty{80}{mL}; \omterm{prop:b2:heart:cycle}{निष्कासन अंश} $80/130 = \qty{62}{\%}$;
निर्गम $80\times 70 = \qty{5.6}{L/min}$। चूकता: \qty{60}{mL},
$60/180 = \qty{33}{\%}$, \qty{4.2}{L/min} --- वह फैला \omterm{def:b2:heart:anatomy}{निलय}
किसी बड़े आयतन पर काम करके निर्गम लगभग बनाए रखता है, पर जितना थामे है
उसका केवल एक तिहाई निकालता है।
\end{solution}

\begin{solution}{exo:b2:heart:6}
$W = 100\times 133\times 7\times 10^{-5} = \qty{0.93}{J}$; शक्ति
$0.93\times 1.2 = \qty{1.1}{W}$। \qty{150}{mmHg} पर: \qty{1.4}{J},
\qty{1.7}{W}। अतिरिक्त \qty{0.56}{W} यांत्रिक यानी \qty{3.7}{W}
उपापचयी, यानी प्रति सेकंड \qty{0.19}{mL} ऑक्सीजन: \qty{11}{mL/min} अधिक।
\end{solution}

\begin{solution}{exo:b2:heart:7}
$60/0.5 = 120$ धड़कनें प्रति मिनट; $60/1.5 = 40$। तीसरा: पूर्ण
हृद्-अवरोध --- \omterm{def:b2:heart:anatomy}{अलिंद} उस गाँठ के नीचे 100 पर धड़कते हैं, \omterm{def:b2:heart:anatomy}{निलय} अपने ही
किसी गतिकारक के नीचे 33 पर, और AV गाँठ कुछ भी चालित नहीं करती।
\end{solution}

\begin{solution}{exo:b2:heart:8}
सरक $20/25 = \qty{0.8}{s}$, धन \qty{0.15}{s}: \qty{0.95}{s}, यानी मिनट
भर में 63 धड़कनें। ऐसिटिलकोलीन: $25/18 = \qty{1.39}{s}$, धन $0.15$:
\qty{1.54}{s}, यानी मिनट भर में 39। नॉरऐड्रिनैलिन: $20/40 = \qty{0.5}{s}$,
धन $0.15$: \qty{0.65}{s}, यानी मिनट भर में 92।
\end{solution}

\begin{solution}{exo:b2:heart:9}
वह देरी \omterm{def:b2:heart:anatomy}{अलिंदों} को \omterm{def:b2:heart:anatomy}{निलयों} के सिकुड़ने से पहले उनमें ख़ाली होना
पूरा करने देती है; उसके बिना \omterm{def:b2:heart:anatomy}{अलिंद} और \omterm{def:b2:heart:anatomy}{निलय} साथ
सिकुड़ते, AV कपाट आधे भरे \omterm{def:b2:heart:anatomy}{निलयों} पर बंद हो जाते
और \omterm{prop:b2:heart:cycle}{आघात-आयतन} अलिंदीय अंशदान भर गिर जाता।
कोई लंबा PR अंतराल केवल निलयीय प्रकुंचन टालता है
और तब तक कुछ नहीं लेता जब तक वह देरी इतनी लंबी न हो जाए कि धड़कनें
गिरने लगें या चालन पूरी तरह चूक जाए।
\end{solution}

\begin{solution}{exo:b2:heart:10}
विश्राम: $5000/45 = \qty{111}{mL}$; व्यायाम: $\num{30000}/180 =
\qty{167}{mL}$। स्टार्लिंग नियम व्यायाम की बड़ी शिरीय वापसी
को किसी बड़े \omterm{prop:b2:heart:cycle}{आघात-आयतन} में बदल देता है; अनुकंपी तंत्रिकाएँ
दर तथा सुकुंचनशीलता (कोई नीचा अंत-प्रकुंचनी आयतन) जोड़ती हैं। लगभग
200 से ऊपर वह अनुशिथिलन \omterm{def:b2:heart:anatomy}{निलय} भरने को बहुत छोटा है और
\omterm{prop:b2:heart:cycle}{आघात-आयतन} उस दर के चढ़ने से तेज़ गिरता है।
\end{solution}

\begin{solution}{exo:b2:heart:11}
$v = Q/A$: $250/3 = \qty{83}{cm/s}$ और $250/1 = \qty{250}{cm/s}$।
$\Delta P = \tfrac12\times 1060\times v^{2}$: \qty{365}{Pa}
(\qty{2.7}{mmHg}) और \qty{3300}{Pa} (\qty{25}{mmHg})। वह \omterm{def:b2:heart:anatomy}{निलय}
वह अतिरिक्त दाब अपनी भित्ति मोटी करके पैदा करता है; और वह मोटा,
कड़ा \omterm{def:b2:heart:anatomy}{निलय} अंततः अपनी हृद्-आपूर्ति से आगे बढ़ जाता है, ख़राब
भरता है, और चूक जाता है।
\end{solution}

\begin{solution}{exo:b2:heart:12}
स्वचालिता बिना किसी निवेश के कोई धड़कन देती है; स्टार्लिंग नियम
निर्गम को वापसी से मिलाता है और दोनों \omterm{def:b2:heart:anatomy}{निलयों} को बिना
किसी निवेश के संतुलित करता है; और तंत्रिकाएँ तथा ऐड्रिनैलिन उस
स्व-नियंत्रित केंद्रक के इर्द-गिर्द केवल दर तथा सुकुंचनशीलता तय करती हैं। कोई प्रतिरोपित
\omterm{def:b2:heart:anatomy}{हृदय} धड़कता है (लगभग 100 पर, बिना वेगसी तान के), अपना
\omterm{prop:b2:heart:cycle}{आघात-आयतन} भरने के अनुसार समंजित करता है, और व्यायाम में अपना निर्गम
स्टार्लिंग क्रियाविधि से तथा परिसंचरित ऐड्रिनैलिन से चढ़ा सकता है --- पर धीरे,
और किसी तंत्रिकित \omterm{def:b2:heart:anatomy}{हृदय} से कम।
\end{solution}

\begin{solution}{pb:b2:heart:1}
\textbf{1.} \qty{70}{mL}; $70/130 = \qty{54}{\%}$; $70\times 70 =
\qty{4.9}{L/min}$।
\textbf{2.} $60/70 = \qty{0.86}{s}$; अनुशिथिलन $0.86 - 0.3 =
\qty{0.56}{s}$।
\textbf{3.} $70/0.3 = \qty{233}{mL/s}$; $233/3 = \qty{78}{cm/s}$।
\textbf{4.} $70/1500 = \qty{0.047}{s}$, यानी लगभग \qty{50}{ms}।
\textbf{5.} $100\times 133\times 7\times 10^{-5} = \qty{0.93}{J}$;
$0.93\times 70/60 = \qty{1.1}{W}$।
\textbf{6.} दोनों \omterm{def:b2:heart:anatomy}{निलय} $1.1\times 1.2 = \qty{1.3}{W}$;
\qty{15}{\%} पर, \qty{8.7}{W} उपापचयी; $8.7/20 = \qty{0.43}{mL}$
ऑक्सीजन प्रति सेकंड, \qty{26}{mL/min}।
\textbf{7.} प्रति मिलिलीटर \omterm{def:b2:blood-circulation:blood}{रक्त} \qty{0.14}{mL} निकालते हुए:
$26/0.14 = \qty{190}{mL/min}$।
\textbf{8.} अंतराल \qty{0.86}{s}, सरक-काल \qty{0.71}{s}:
$20/0.71 = \qty{28}{mV/s}$।
\textbf{9.} अंतराल \qty{0.6}{s}, सरक \qty{0.45}{s}: \qty{44}{mV/s}।
\textbf{10.} अंतराल \qty{0.33}{s}, सरक \qty{0.18}{s}: \qty{110}{mV/s}।
\textbf{11.} R--R \qty{0.86}{s}; P: अलिंदीय विध्रुवण; QRS:
निलयीय विध्रुवण; T: निलयीय पुनर्ध्रुवण।
\textbf{12.} वह देरी \omterm{def:b2:heart:anatomy}{अलिंदों} को \omterm{def:b2:heart:anatomy}{निलयों} के सिकुड़ने से पहले
उनमें ख़ाली होने देती है; 180 पर पूरा चक्र \qty{0.33}{s} है, इसलिए कोई
अपरिवर्तित \qty{0.16}{s} देरी उसका आधा खा जाती --- इसीलिए
अनुकंपी तंत्रिकाएँ उस गाँठ का चालन भी तेज़ कर देती हैं।
\textbf{13.} पूर्ण हृद्-अवरोध: \omterm{def:b2:heart:anatomy}{निलय} मिनट भर में 40 पर,
बंडल या \omterm{prop:b2:heart:conduction}{पुर्किंजे तंतुओं} से गति पाते, और \omterm{def:b2:heart:anatomy}{अलिंद}
\omterm{prop:b2:heart:conduction}{साइनु-अलिंद गाँठ} के नीचे 75 पर।
\textbf{14.} \qty{120}{mL}; $120/150 = \qty{80}{\%}$; $120\times 180 =
\qty{21.6}{L/min}$।
\textbf{15.} $0.33 - 0.3 = \qty{0.03}{s}$: भरने को लगभग कोई समय नहीं;
और तेज़ हो तो \omterm{prop:b2:heart:cycle}{आघात-आयतन} ढह जाता है।
\textbf{16.} $W = 120\times 133\times 1.2\times 10^{-4} = \qty{1.9}{J}$,
प्रति सेकंड $\times 3$: \qty{5.7}{W}; दोनों \omterm{def:b2:heart:anatomy}{निलय} \qty{6.9}{W};
उपापचयी \qty{46}{W}; ऑक्सीजन \qty{2.3}{mL/s}, \qty{140}{mL/min}।
\textbf{17.} $140/0.14 = \qty{1000}{mL/min}$, यानी विश्राम-प्रवाह से
पाँच गुना, और वह भी चक्र के दसवें भाग तक सिकुड़े किसी अनुशिथिलन में देना है:
हृद्-धमनिकाओं को अपनी सीमा तक फैलना पड़ेगा।
\textbf{18.} बिना वेगसी तान के दर लगभग 100, अंत-प्रकुंचनी आयतन
\qty{60}{mL} पर अपरिवर्तित: \omterm{prop:b2:heart:cycle}{आघात-आयतन} $150 - 60 = \qty{90}{mL}$,
निर्गम \qty{9}{L/min}।
\textbf{19.} केवल दर (\qty{70}{mL} पर $70 \to 180$, यानी
\qty{12.6}{L/min}): $+7.7$; भरना (180 पर \qty{20}{mL} अधिक):
$+3.6$; सुकुंचनशीलता (180 पर \qty{30}{mL} कम अवशेष): $+5.4$ ---
यानी 4.9 से \qty{21.6}{L/min}, और तीनों अंशदान ठीक
$+16.7$ जुड़ते हैं।
\textbf{20.} $233/1 = \qty{233}{cm/s}$, \qty{2.3}{m/s}।
\textbf{21.} $\tfrac12\times 1060\times 2.33^{2} = \qty{2900}{Pa}$,
\qty{22}{mmHg}।
\textbf{22.} $120/0.3 = \qty{400}{mL/s}$, \qty{4}{m/s}; $\tfrac12\times
1060\times 16 = \qty{8500}{Pa}$, \qty{64}{mmHg}।
\textbf{23.} शिखर विश्राम में $120 + 22 = \qty{142}{mmHg}$, व्यायाम में लगभग $140 +
64 = \qty{204}{mmHg}$; विश्राम में माध्य निष्कासन दाब
100 से चढ़कर 122: प्रति धड़कन काम $\times 1.22$।
\textbf{24.} कोई मोटी भित्ति अधिक बल पैदा करती है और प्रति तंतु
प्रतिबल घटाती है; पर सींचा जाने वाला पिंड बढ़ता है जबकि
हृद्-प्रवाह, जो प्रकुंचन में उसी मोटी भित्ति से भिंचता है और जिसे कोई
छोटा अनुशिथिलन मिलता है, साथ नहीं चल पाता, और कोई मोटी भित्ति ख़राब
ढीली होती और कम भरती है: रक्ताभाव तथा विफलता आती है।
\textbf{25.} विश्राम में \qty{4.9}{L/min}, व्यायाम में \qty{21.6}{L/min};
प्रति धड़कन \qty{0.93}{J}; प्रवणता विश्राम में \qty{22}{mmHg} और
व्यायाम में \qty{64}{mmHg}।
\end{solution}
